% Copyright (c) 2026 Benjamin Stanley Frohman.
% Licensed under the Apache License, Version 2.0.
% REVIEW DRAFT --- not a DOI assignment, upload, or publication.
\documentclass[11pt,a4paper]{article}

\usepackage[margin=1.05in]{geometry}
\usepackage{amsmath,amssymb,amsthm,mathtools}
\usepackage{microtype}
\usepackage{enumitem}
\usepackage{booktabs}
\usepackage{longtable}
\usepackage{hyperref}
\usepackage[T1]{fontenc}
\usepackage{lmodern}
\usepackage{setspace}
\usepackage{xcolor}
\usepackage{xurl}
\Urlmuskip=0mu plus 3mu
\setstretch{1.08}

\hypersetup{colorlinks=true,linkcolor=black,citecolor=black,urlcolor=blue}

\newcommand{\Q}{\mathbb{Q}}
\newcommand{\C}{\mathbb{C}}
\newcommand{\lean}[1]{\texttt{\detokenize{#1}}}

\newtheorem{theorem}{Theorem}
\newtheorem{lemma}[theorem]{Lemma}
\newtheorem{corollary}[theorem]{Corollary}
\theoremstyle{definition}
\newtheorem{definition}[theorem]{Definition}
\newtheorem{remark}[theorem]{Remark}

\title{\textbf{Exact B\'ezout Certificates and Coefficient Nonvanishing\\
at Section Roots}\\[0.35em]
\large Lean Stage~2 Algebraic Certificates\\[0.4em]
\normalsize\textcolor{red}{\textbf{REVIEW DRAFT}} --- not peer-reviewed; no DOI assigned; not uploaded}
\author{Frohman, Benjamin Stanley\\
\small ORCID: \href{https://orcid.org/0009-0006-7068-3718}{0009-0006-7068-3718}\\
\small Email: \texttt{frohmanbenjamin@gmail.com}\\
\small Independent Researcher}
\date{6~November~2026 (America/Chicago)\\
\small Source pin: \texttt{e207394a9ed4255ad804860b8f69fde7e940dd54}\\
\small Repository: \url{https://github.com/BenFrohman/HODGE}\\
\small Pull request: \url{https://github.com/BenFrohman/HODGE/pull/36}}

\begin{document}
\maketitle

\begin{center}
\fbox{\parbox{0.92\textwidth}{\centering\small
\textbf{Status.} This manuscript is a \textbf{REVIEW DRAFT} prepared for later
integration with Zenodo packaging. It invents no DOI, performs no upload or
publication, and does not claim peer review, novelty, or a Clay/Hodge close.
License of the cited Lean sources remains Apache-2.0; this note does not
relicense them. Author metadata uses the exact form
\texttt{Frohman, Benjamin Stanley}.}}
\end{center}

\begin{abstract}
We record the Stage~2 Lean formalization of exact rational polynomial
reduction and B\'ezout identities for nine explicit section-coefficient
polynomials associated to a fixed degree-$24$ section modulus, together with
the resulting nonvanishing theorems at every root of that modulus, the proof
that the modulus has degree~$24$, the existence of a complex root, and the
existence of a complex root at which all nine quantities are simultaneously
nonzero. The module
\path{Hodge.ExplicitResidueCertificates} contains $45$ theorem declarations;
combined with the seven Stage~1 lemmas this yields the audited total of $52$
incremental theorems, each depending only on the standard Lean axioms
\lean{propext}, \lean{Classical.choice}, and \lean{Quot.sound}. SymPy is used
only to generate candidate witnesses; Lean proves every identity
independently. No geometric cycle-class, Jacobian-quotient basis, or H/G
resultant theorem is claimed.
\end{abstract}

\tableofcontents

\section{Scope and non-claims}

This note is the companion to the Stage~1 determinant/derivative note. It
covers every new Stage~2 theorem in
\lean{Hodge/ExplicitResidueCertificates.lean} at commit
\path{e207394a9ed4255ad804860b8f69fde7e940dd54}, organized as one coherent
algebraic certificate paper. Generated helper declarations (evaluation lemmas,
reduction/B\'ezout pairs) are mapped as supporting infrastructure for the nine
nonvanishing results and three nonvacuity theorems; they are not split into
separate papers~\cite{HODGERepo}.

\paragraph{Proved here.}
\begin{itemize}[leftmargin=1.4em,itemsep=0.25em]
\item Nine scalar nonvanishing theorems for the formulas $w_0,w_1,w_2,w_3$,
      $u_0,u_1,u_2,u_3$, and $\mathrm{disc}=b^2-4c$, under only the field /
      $\Q$-algebra structure and the modulus-root equation.
\item Eighteen exact identities in $\Q[X]$: for each of the nine raw
      polynomials $P$, a reduction $P=R+Q\cdot m$ and a B\'ezout identity
      $A R+B m=1$.
\item Degree $24$ of the modulus polynomial, existence of a complex root, and
      simultaneous nonvanishing of all nine quantities at some complex root.
\end{itemize}

\paragraph{Not proved here.}
\begin{itemize}[leftmargin=1.4em,itemsep=0.25em]
\item The tenth Python certificate: the actual H/G resultant and its
      common-projective-zero consequence. Pinned mathlib defines resultant via
      a Sylvester determinant and lists further resultant theory as TODO
      (\lean{docs/FULL_FORMALIZATION_STAGE2.md}).
\item Identification of the $U$ formulas with coefficients in an actual
      Jacobian quotient (requires a Lean quotient-reduction/basis theorem).
\item Complete-intersection ideal regularity, reducedness, saturation;
      geometric residue-cycle theorems; Betti/Hodge comparison; geometric group
      action; rational geometric spanning of Hodge classes.
\item General or integral Hodge conjecture statements; novelty or peer-review
      claims.
\end{itemize}

\section{Standing data}

\begin{definition}[Modulus and section coefficients]
Let $m\in\Q[X]$ be the explicit polynomial
\begin{equation}
m = X^{24}+66081312000\,X^{18}-4811512860000\,X^{12}
+1167566400000\,X^{6}+583200000
\end{equation}
(\lean{ExplicitResidueCertificates.modulus}). Let
$b,c,d,e,g\in\Q[X]$ be the section polynomials
\lean{bPoly},\lean{cPoly},\lean{dPoly},\lean{ePoly},\lean{gPoly}, whose
evaluations recover \lean{ExplicitCurve.sb},\ldots,\lean{sg}. For a
characteristic-zero field $K$ equipped with a $\Q$-algebra structure
and $k\in K$, write
$b(k)=\mathrm{sb}(k)$ etc., and define the nine scalar formulas
\begin{align}
w_0 &= 3bg-2ce, &
w_1 &= 6g+be-4cd, &
w_2 &= 4e-bd-6c, &
w_3 &= 2d-3b,\\
u_0 &= c w_0-6, &
u_1 &= b w_0+c w_1, &
u_2 &= w_0+b w_1+c w_2, &
u_3 &= w_1+b w_2+c k,\\
\mathrm{disc} &= b^2-4c. &&&&
\end{align}
The corresponding elements of $\Q[X]$ are \lean{w0Poly},\ldots,\lean{discPoly}.
\end{definition}

\begin{remark}[Trust boundary for SymPy]
The script \path{scripts/generate_residue_certificates.py} uses
SymPy~$1.14.0$~\cite{SymPy} to compute candidate sparse rational
polynomials $R,Q,A,B$. Those candidates are written into Lean source;
every equality is then proved by Lean polynomial extensionality and exact
ring normalization. SymPy is reproducible generation tooling, not part of
the proof trust base (\path{docs/FULL_FORMALIZATION_STAGE2.md}).
\end{remark}

\section{Supporting polynomial identities}

\begin{lemma}[Evaluation]\label{lem:eval}
For each section polynomial $P\in\{b,c,d,e,g\}$ and each raw coefficient
polynomial $P\in\{w_0,\ldots,w_3,u_0,\ldots,u_3,\mathrm{disc}\}$,
\begin{equation}
\mathrm{aeval}_k(P_{\mathrm{Poly}})=P(k)
\end{equation}
in every characteristic-zero field $K$ equipped with a $\Q$-algebra
structure. {\sloppy Lean declares the hypotheses
\texttt{[Field K]}, \texttt{[CharZero K]}, and
\texttt{[Algebra}~$\mathbb{Q}$~\texttt{K]}.}
In particular the five section polynomials evaluate to
\lean{ExplicitCurve.sb},\ldots,\lean{sg}.
\end{lemma}
\begin{proof}[Lean]
Fourteen \lean{@[simp]} theorems
(\lean{eval_bPoly} through \lean{eval_disc})
are proved under those field and $\Q$-algebra hypotheses by unfolding
definitions and applying the algebra-map simp set.
\end{proof}

\begin{lemma}[Reduction and B\'ezout]\label{lem:bezout}
For each of the nine raw polynomials $P\in\Q[X]$ there exist explicitly
defined sparse $R,Q,A,B\in\Q[X]$ such that
\begin{equation}
P = R + Q\cdot m
\qquad\text{and}\qquad
A R + B m = 1
\end{equation}
as equalities in $\Q[X]$.
\end{lemma}
\begin{proof}[Lean]
Eighteen theorems
\path{w0_reduction}, \path{w0_bezout}, \ldots,
\path{disc_reduction}, \path{disc_bezout}
are proved by \lean{Polynomial.funext}, evaluation at a generic point,
and \lean{norm_num}/\lean{ring}. Exact rational coefficients are recorded
in the Lean source; they are omitted here only for readability.
\end{proof}

\begin{lemma}[Evaluation criterion]\label{lem:criterion}
If $P=R+Q m$ and $A R+B m=1$ in $\Q[X]$, and if $k$ is a root of $m$ in a
characteristic-zero field $K$ equipped with a $\Q$-algebra structure
(Lean: \texttt{[Field K]}, \texttt{[CharZero K]},
\texttt{[Algebra}~$\mathbb{Q}$~\texttt{K]}),
then $\mathrm{aeval}_k(P)\neq 0$.
\end{lemma}
\begin{proof}[Lean]
\lean{eval_ne_zero_of_certificate}, under the same field/$\Q$-algebra
hypotheses: evaluate both identities at $k$; the B\'ezout identity
collapses to $A(k)R(k)=1$, so $R(k)\neq0$, hence $P(k)\neq0$.
\end{proof}

\section{Main nonvanishing theorems}

\begin{theorem}[Nine coefficient nonvanishings]\label{thm:nine}
Let $K$ be a field of characteristic zero with a $\Q$-algebra structure, and
let $k\in K$ satisfy
\begin{equation}
k^{24}+66081312000\,k^{18}-4811512860000\,k^{12}
+1167566400000\,k^{6}+583200000=0.
\end{equation}
Then
\begin{equation}
w_0(k),\;w_1(k),\;w_2(k),\;w_3(k),\;
u_0(k),\;u_1(k),\;u_2(k),\;u_3(k),\;
\mathrm{disc}(k)
\end{equation}
are all nonzero.
\end{theorem}
\begin{proof}[Lean]
Nine theorems \lean{w0_ne_zero},\ldots,\lean{disc_ne_zero}. Each rewrites the
scalar hypothesis as $\mathrm{aeval}_k(m)=0$ and applies
Lemma~\ref{lem:criterion} to the matching reduction/B\'ezout pair. No
numerical root approximation, irreducibility hypothesis, geometric datum, or
desired nonvanishing is assumed.
\end{proof}

\begin{theorem}[Modulus degree]\label{thm:degree}
$\mathrm{degree}(m)=24$.
\end{theorem}
\begin{proof}[Lean]
\lean{modulus_degree}: unfold $m$, apply \lean{compute_degree}, and discharge
closed numerical bounds by \lean{norm_num} and ordinary
\lean{decide +kernel}. This is not a \lean{native_decide} shortcut on the
algebraic content of the nonvanishing theorems.
\end{proof}

\begin{theorem}[Complex root existence]\label{thm:exists}
There exists $k\in\C$ with $\mathrm{aeval}_k(m)=0$.
\end{theorem}
\begin{proof}[Lean]
\lean{exists_complex_root}: apply algebraic closedness of $\C$
(\path{IsAlgClosed.exists_aeval_eq_zero}) using Theorem~\ref{thm:degree} and
$\mathrm{degree}\,m\neq0$.
\end{proof}

\begin{theorem}[Simultaneous nonvanishing at a complex root]\label{thm:simul}
There exists $k\in\C$ satisfying the modulus equation such that
\begin{equation}
w_0(k)\neq0,\;\ldots,\;w_3(k)\neq0,\;
u_0(k)\neq0,\;\ldots,\;u_3(k)\neq0,\;
\mathrm{disc}(k)\neq0.
\end{equation}
\end{theorem}
\begin{proof}[Lean]
\lean{exists_complex_nonvanishing}: take a root from
Theorem~\ref{thm:exists} and apply all nine cases of
Theorem~\ref{thm:nine}. This shows the universal nonvanishing statements are
not vacuous assertions over a field without section roots. It is not a
geometric cycle-class assertion.
\end{proof}

\section{Theorem--declaration mapping}
\label{sec:mapping}

\subsection*{Main results}

\begin{center}
\begin{longtable}{@{}llp{6.2cm}@{}}
\toprule
\# & Lean declaration & Role \\
\midrule
1 & \lean{w0_ne_zero} & $w_0=3bg-2ce\neq0$ \\
2 & \lean{w1_ne_zero} & $w_1=6g+be-4cd\neq0$ \\
3 & \lean{w2_ne_zero} & $w_2=4e-bd-6c\neq0$ \\
4 & \lean{w3_ne_zero} & $w_3=2d-3b\neq0$ \\
5 & \lean{u0_ne_zero} & $u_0=cw_0-6\neq0$ \\
6 & \lean{u1_ne_zero} & $u_1=bw_0+cw_1\neq0$ \\
7 & \lean{u2_ne_zero} & $u_2=w_0+bw_1+cw_2\neq0$ \\
8 & \lean{u3_ne_zero} & $u_3=w_1+bw_2+ck\neq0$ \\
9 & \lean{disc_ne_zero} & $b^2-4c\neq0$ \\
10 & \lean{modulus_degree} & $\deg m=24$ \\
11 & \lean{exists_complex_root} & $\exists\,k\in\C$, $m(k)=0$ \\
12 & \lean{exists_complex_nonvanishing} & simultaneous nonvanishing \\
\bottomrule
\end{longtable}
\end{center}

\subsection*{Supporting certificates (same paper; not split)}

\begin{center}
\begin{longtable}{@{}llp{5.8cm}@{}}
\toprule
Count & Lean declarations & Role \\
\midrule
5 & \lean{eval_bPoly}\ldots\lean{eval_gPoly} & section evaluation \\
9 & \lean{eval_w0}\ldots\lean{eval_disc} & raw-coefficient evaluation \\
9 & \lean{*_reduction} & $P=R+Qm$ in $\Q[X]$ \\
9 & \lean{*_bezout} & $AR+Bm=1$ in $\Q[X]$ \\
1 & \lean{eval_ne_zero_of_certificate} & evaluation criterion \\
\bottomrule
\end{longtable}
\end{center}

\noindent
Total Stage~2 theorems: $12+5+9+9+9+1=45$. Combined with the seven Stage~1
lemmas: $52$ audited declarations
(\lean{ResidueAxiomAudit.lean}).

\section{Axiom audit and verification evidence}
{\sloppy

The checker \path{scripts/check_residue_axioms.py} requires
\lean{ResidueAxiomAudit.lean} to enumerate every theorem in both incremental
modules in source order, builds
\path{Hodge.ExplicitResidueAlgebra} and
\path{Hodge.ExplicitResidueCertificates}, and accepts only
\lean{propext}, \lean{Classical.choice}, \lean{Quot.sound}.

Verified log lines (America/Chicago evidence timestamps on disk):
\begin{itemize}[leftmargin=1.4em]
\item \path{build-and-audit-final.log}: reports successful
      \path{Hodge.ExplicitResidueCertificates} build ($7085/7085$) and
      ``PASS: 52 incremental theorem declarations; standard axioms only''.
\item \path{residue-audit-e207394.log}: same PASS line; $52$ axiom-print
      entries.
\item \lean{no-sorry-e207394.log}:
      ``OK: no sorry/admit/native\_decide, no axiom declarations''.
\item \lean{full-build-e207394.log}: full repository
      \lean{lake build} completed successfully at commit \lean{e207394}.
\item Prior regeneration check (certificate delivery evidence):
      byte-identical SymPy regeneration of
      \lean{ExplicitResidueCertificates.lean}; \lean{LICENSE},
      \lean{CITATION.cff}, and Stage~1 Lean source unchanged.
\end{itemize}

} % end sloppy

\section{Toolchain pins and reproducibility}
{\sloppy

\begin{center}
\begin{tabular}{@{}ll@{}}
\toprule
Pin & Value \\
\midrule
Lean~\cite{Lean4} & \lean{leanprover/lean4:v4.22.0} \\
mathlib input~\cite{mathlib4} & \lean{v4.22.0} \\
mathlib checkout & \path{79e94a093aff4a60fb1b1f92d9681e407124c2ca} \\
SymPy (generator only)~\cite{SymPy} & $1.14.0$ \\
Git commit & \path{e207394a9ed4255ad804860b8f69fde7e940dd54} \\
PR & \url{https://github.com/BenFrohman/HODGE/pull/36} \\
License & Apache-2.0 (preserved; no relicense) \\
\bottomrule
\end{tabular}
\end{center}

\begin{verbatim}
python3 scripts/generate_residue_certificates.py
# optional: git diff --exit-code -- Hodge/ExplicitResidueCertificates.lean
lake build Hodge.ExplicitResidueAlgebra \
           Hodge.ExplicitResidueCertificates
python3 scripts/check_residue_axioms.py
python3 scripts/check_no_sorry.py
\end{verbatim}

Local sources and evidence inspected (read-only; no repository edits).
Mathematical source pin \path{e207394a9ed4255ad804860b8f69fde7e940dd54};
later inspected worktree \path{fefe2ba9d87c900b538d02c13f7f51f375fcd285}
retains identical reviewed modules, audit, and toolchain/lock files:
\begin{itemize}[leftmargin=1.4em]
\item \path{/workspace/hodge-formal-stage1-20261106/Hodge/ExplicitResidueCertificates.lean}
\item \path{/workspace/hodge-formal-stage1-20261106/ResidueAxiomAudit.lean}
\item \path{/workspace/hodge-formal-stage1-20261106/docs/FULL_FORMALIZATION_STAGE2.md}
\item \path{/workspace/hodge-formal-stage1-20261106/scripts/generate_residue_certificates.py}
\item \path{/workspace/hodge-formal-stage2-evidence-20261106/}
      (actual files: \lean{build-and-audit-final.log},
      \lean{build-and-audit.log}, \lean{full-build.log},
      \lean{typecheck*.log}, \lean{DegreeCheck.lean},
      \lean{generate_certificates.py})
\item \path{/workspace/hodge-formal-stage1-evidence-20261106/}
      (\lean{Stage1Audit.lean}, \lean{axioms.txt},
      \lean{baseline-explicit-curve.lean})
\item \path{/workspace/hodge-certificate-delivery-20261106/}
      (\lean{full-build-e207394.log},
      \lean{residue-audit-e207394.log},
      \lean{no-sorry-e207394.log}, status JSON).
\end{itemize}

} % end sloppy

\paragraph{CI verification at a later identical head.}
User-reported merge of PR~\url{https://github.com/BenFrohman/HODGE/pull/36}
is acknowledged; this packaging records no merge-commit SHA (merge
verification is underway separately). At head
\path{fefe2ba9d87c900b538d02c13f7f51f375fcd285}, where the reviewed Lean
modules, audit, and toolchain/lock match mathematical pin
\path{e207394a9ed4255ad804860b8f69fde7e940dd54}, all five reported GitHub
Actions checks completed successfully:
Lean/build and Lean/attempt
(\url{https://github.com/BenFrohman/HODGE/actions/runs/42317080169}),
Lean certificate/certify on the PR
(\url{https://github.com/BenFrohman/HODGE/actions/runs/42317080616}),
Lean certificate/certify on the push
(\url{https://github.com/BenFrohman/HODGE/actions/runs/42317075046}),
and Paper/paper
(\url{https://github.com/BenFrohman/HODGE/actions/runs/42317080271}).
Both certification logs build the incremental modules and print
\texttt{PASS: 52 incremental theorem declarations; standard axioms only}.
CI success verifies those algebraic declarations; it does not discharge
geometric Hodge/Betti/cycle-class obligations. This note remains a
REVIEW DRAFT with no DOI and no publication claim.

\section{Limitations relative to the geometric program}

The nine nonvanishings are algebraic statements about evaluations of
explicit elements of $\Q[X]$ at roots of $m$. They do not by themselves
produce geometric residues, nor do they prove that the corresponding
Jacobian-quotient classes are nonzero in a geometric cohomology theory.
Villaflor Loyola's complete-intersection residue theorem~\cite{Villaflor2021}
and Griffiths' comparison remain external classical input
(\lean{NOTICE.md}).

The classical Python projection certificate checks ten exact unit-at-every-root
identities (eight $W$/$U$ formulas, the discriminant, and the H/G resultant).
Lean Stage~2 covers the first nine; the tenth (resultant) remains unformalized.

These algebraic certificates can stand as fully formalized Lean mathematics
while a full geometric Hodge paper for the sextic fourfold remains unpublished
pending the missing geometric foundations. This REVIEW DRAFT assigns no DOI
and makes no upload.

\section*{Author, license, and contact}
Author metadata: \textbf{Frohman, Benjamin Stanley}.
Email: \texttt{frohmanbenjamin@gmail.com}.
ORCID: \href{https://orcid.org/0009-0006-7068-3718}{0009-0006-7068-3718}.
Verified from repository \path{CITATION.cff}, \path{NOTICE}, and
\path{.zenodo.json}.
Copyright (c) 2026 Benjamin Stanley Frohman. Apache License 2.0.
Historical archive DOI \texttt{10.5281/zenodo.23629480} is provenance of
prior HODGE releases, not a DOI for this Stage~2 note.

\bibliographystyle{plain}
\bibliography{references}

\end{document}
